\section{Discussion}
\label{sec:discussion}

The principles derived in the preceding sections originate from the development and
investigation of a controlled transformation architecture for heterogeneous legacy
engineering information into SysML v2. Their relevance beyond this specific application
therefore needs to be distinguished from those aspects that remain bound to the investigated
engineering task.

\subsection{Transferability and Context Dependence}
\label{sec:discussion_transferability}

Several mechanisms investigated in the underlying architecture are not inherently tied to
SysML v2 as a target representation. Source processing, provenance, stable identities,
persisted processing states, probabilistic interpretation, and explicit human authorization
address more general problems that also occur in other information-centered engineering
tasks.

This provides a basis for considering similar mechanisms in workflows such as requirements
review, technical-documentation analysis, or quality and regulatory documentation review.
Such tasks likewise involve heterogeneous documents that have to be interpreted in relation
to an engineering context and whose resulting assessments may require explicit human
decisions.

The transferability of these mechanisms should nevertheless be understood as a further
hypothesis rather than as an empirically demonstrated property across engineering domains.
The underlying investigation was limited to a specific brownfield transformation task and
to a defined SysML-v2 target.

In particular, task-specific aspects cannot simply be transferred unchanged. The relevant
Engineering Sources and Context Sources, the perspectives used for probabilistic
interpretation, the authority boundaries, the controlled engineering state, the
target-specific rules, and the assessment mechanisms all depend on the engineering
purpose.

The resulting distinction is important. A potentially transferable aspect is not one fixed
workflow configuration, but the separation of functions: controlled source processing,
probabilistic interpretation, explicit engineering decisions, persistent states,
rule-bound downstream processing, provenance, traceability, and separate assessment.

\subsection{From Controlled Prototype to Productive Engineering Use}
\label{sec:discussion_productive_use}

The investigated architecture demonstrates a controlled processing path under bounded
conditions. Productive use introduces additional questions that were only addressed to a
limited extent in the underlying work.

A first challenge concerns the information basis itself. Productive environments contain
larger and less curated information sets, additional document types, changing information
states, and potentially broader sets of engineering constructs. The mechanisms for source
classification, interpretation, authorization, and controlled integration therefore need to
remain effective as the amount and diversity of information increase.

A second challenge concerns human review. Explicit authorization provides a control
boundary, but it also creates review effort. As the processing volume increases, the
scalability of human review and authorization becomes part of the engineering problem.
The relevant question is therefore not whether the human expert should be removed from the
process, but how human attention can be focused on those states in which engineering
judgement is actually required.

A third challenge is the integration into an existing engineering workflow. Domain
references, interpretation perspectives, rules, representation frameworks, review
responsibilities, and release mechanisms need to fit the respective organizational
context. The practical quality of such a workflow therefore depends not only on model
performance but also on the comprehensibility of the provided decision basis, the actual
review effort, and the compatibility of the authority boundaries with established working
practices. Prior work on MBSE adoption likewise highlights the relevance of organizational
context, perceived value, and adoption conditions when engineering methods are introduced
into practice \cite{Henderson2024,Call2022,Campo2023,Nanfuka2023}.

Finally, productive use also requires the conventional qualities expected from an
engineering software system. Software quality, runtime behavior, resource demand,
robustness, maintainability, and operational integration remain relevant independently of
the AI-specific architecture.

These aspects reinforce the main observation of this paper: moving from an AI-supported
demonstrator to productive engineering use is not merely a matter of improving the
underlying model. The information basis, decision boundaries, engineering rules,
assessment mechanisms, and organizational workflow have to evolve together.

\subsection{Learning from Human Engineering Decisions}
\label{sec:discussion_learning}

Persisted human decisions create an additional long-term perspective. Over repeated use,
an AI-supported engineering workflow can accumulate a structured data basis that connects
engineering information, machine-generated proposals, and the corresponding human
evaluations or corrections.

Research on learning from human feedback demonstrates that human evaluations can provide
useful learning signals for adaptive systems \cite{Christiano2017,Ouyang2022}. In an
engineering context, however, such feedback remains meaningful only if it can be related to
the information state and decision context from which it originated.

The provenance and state persistence required for traceability can therefore also provide
a basis for controlled learning. Repeated corrections may indicate that interpretation
instructions, examples, rules, or specialized models should be adapted. Such adaptations
should themselves be versioned, evaluated separately, and explicitly released before they
become part of the productive workflow.

With increasing quality of machine-generated interpretations, human review may eventually
be concentrated more strongly on cases characterized by relevant uncertainty, variance,
contradictory information, or novel situations. Whether individual decisions can be further
automated should, however, be determined empirically from future application data rather
than assumed from model capability alone.

This perspective extends the architecture from a controlled transformation process toward
a potentially adaptive engineering workflow while retaining the same fundamental
requirement: changes in automated behavior must remain attributable, assessable, and
subject to explicit engineering control.
